%========================================================
% CHAPTER 2 — Modeling Strategy
%========================================================

\chapter{Modeling Strategy}

%========================================================
% SECTION 2.1 — Progressive Modeling Philosophy
%========================================================

\section{Progressive Modeling Philosophy}

The modeling approach adopted in this work follows a progressive refinement strategy. The physical system is not introduced directly through the most complete effective description. Instead, the fiber model is constructed by adding one propagation mechanism at a time, so that the role of each contribution can be isolated analytically and validated numerically.

\medskip

The central physical object is an entangled photon pair distributed through two independent optical-fiber arms. The fiber action is described at the level of the polarization qubits. Within this setting, the relevant effects are organized into coherent phase evolution, statistical phase fluctuations, and depolarizing noise.

\medskip

This hierarchy serves two purposes. From the theoretical point of view, it separates distinct physical mechanisms and avoids mixing coherent, statistical, and noisy effects at the same modeling level. From the computational point of view, it leads to a modular simulator architecture, where each new model extends a previously validated implementation.

\medskip

The first levels of the hierarchy provide analytically tractable reference cases. Later levels introduce effective noise mechanisms and composite fiber-channel descriptions, progressively improving the representation of two photons propagating through independent optical fibers.

%========================================================
% SECTION 2.2 — Modeling Hierarchy
%========================================================

\section{Modeling Hierarchy}

The starting point is the ideal Bell-state preparation. This defines the reference entangled state against which all propagation and degradation effects are evaluated.

\medskip

The model is then refined through successive levels:

\[
\begin{array}{c}
\text{Ideal Bell-state preparation} \\
\downarrow \\
\text{Deterministic reduced phase evolution} \\
\downarrow \\
\text{Random phase ensemble} \\
\downarrow \\
\text{Global depolarizing baseline} \\
\downarrow \\
\text{Local two-arm depolarizing channel} \\
\downarrow \\
\text{Composite effective fiber channel}
\end{array}
\]

\medskip

The first propagation model is the deterministic reduced phase model. After polarization compensation, the dominant residual coherent effect is represented as a relative phase shift between the Bell-state components. This model preserves the purity of the global state and isolates the phase dependence of the polarization correlations.

\medskip

The next refinement introduces statistical phase fluctuations. Instead of assigning a single deterministic value to the phase parameter, the state is averaged over a distribution of sampled phase realizations. This leads naturally to a density-operator description and models loss of coherence due to unresolved phase variability.

\medskip

Depolarizing noise is then introduced as an effective description of polarization degradation. The global depolarizing channel is used as an analytically simple isotropic-noise baseline, while the local two-arm depolarizing channel acts independently on arms \( A \) and \( B \), reflecting the physical geometry of two photons propagating through distinct fibers.

\medskip

The final developed model is the composite effective fiber channel. It combines residual phase evolution, local two-arm depolarization, and, when required, statistical phase averaging. This model represents the main effective fiber-propagation description considered in the numerical experiments and leaves the framework open to future extensions.

%========================================================
% SECTION 2.3 — From Model Levels to Numerical Experiments
%========================================================

\section{From Model Levels to Numerical Experiments}

Each model level is associated with a dedicated numerical experiment. The experiments are not independent demonstrations, but successive validation stages of the same modeling framework.

\medskip

This organization makes the numerical part of the thesis cumulative. Each experiment introduces a single additional modeling ingredient with respect to the previous ones, while preserving the same computational structure: state preparation, evolution, observable evaluation, metric computation, and comparison with analytical predictions.

\medskip

\begin{table}[htbp]
\centering
\small
\setlength{\tabcolsep}{4pt}
\begin{tabular}{lll}
\hline
\textbf{Exp.} & \textbf{Model level} & \textbf{New insertion} \\
\hline
1 & Reduced phase model & Deterministic phase sweep \\
2 & Random phase ensemble & Statistical averaging over \( \theta_k \) \\
3 & Global depolarization & Isotropic mixed-state baseline \\
4 & CHSH analysis & Bell-inequality diagnostics from \( T \) \\
5 & Composite fiber channel & Phase plus local two-arm depolarization \\
\hline
\end{tabular}
\caption{Connection between the modeling hierarchy and the numerical experiments. Each experiment introduces only the additional mechanism required at that stage.}
\label{tab:modeling_hierarchy_experiments}
\end{table}

\medskip

The first experiments establish reference cases in which analytical predictions can be derived explicitly and used to validate the simulator. Later experiments reuse the same validated components in more structured models, where coherent phase evolution, statistical averaging, and effective noise are combined.

\medskip

This structure ensures that each numerical result can be interpreted as the consequence of a specific modeling refinement. The analytical predictions developed in the theoretical chapters are therefore tested progressively before being combined in the composite fiber-channel model.